% !TEX root = relatorio.tex
%
% Capítulo 8 — Conclusão
%
\chapter{Conclusão} \label{cap:conclusao}

O Play4Change nasceu de uma lacuna concreta: não existia, até ao momento, uma plataforma de aprendizagem adaptativa que integrasse nativamente o referencial DigComp~2.2 e recorresse à Inteligência Artificial para gerar e personalizar conteúdo a partir de fontes documentais arbitrárias, em vez de um catálogo estático curado por especialistas. O projeto respondeu a essa lacuna com uma plataforma multi-cliente completa --- painel de administração web, aplicação móvel para Learners e um \textit{backend} que orquestra ingestão de conteúdo, geração por IA, indexação vetorial e um percurso de remediação adaptativa --- e validou-a tecnicamente de ponta a ponta, da arquitetura ao utilizador final.

Os cinco objetivos operacionais definidos no Capítulo~\ref{cap:intro} foram todos alcançados, com o detalhe da concretização de cada um analisado no Capítulo~\ref{cap:discussao} (§\ref{sec:analise-objetivos}): a plataforma adaptativa está operacional e alinhada com o DigComp~2.2 ao nível da curadoria de tópicos, confirmada por uma avaliação de usabilidade em que 66,7\% dos participantes classificaram o sistema como \textit{Bom} ou \textit{Excelente} (SUS, mediana~85,0); o \textit{pipeline} de geração de conteúdo por IA funciona de ponta a ponta sobre fontes arbitrárias --- PDF ou URL --- com seis garantias estruturais de qualidade por questão gerada; a remediação adaptativa classifica o padrão de erro do Learner e gera sub-tarefas dirigidas a esse padrão; a tutoria conversacional por IA escala a remediação quando esta se esgota, com isolamento tipado do \textit{input} do utilizador; e a plataforma está protegida por um modelo de segurança em profundidade, validado por uma análise de risco OWASP que não deixou nenhum risco residual em nível Alto ou Crítico.

Do ponto de vista técnico, três decisões distinguem-se como as mais determinantes para o resultado final. A arquitetura hexagonal, com o domínio isolado de qualquer \textit{framework} e o contrato \texttt{Either<DomainError,~T>} da Arrow a tornar os casos de insucesso tão explícitos como os de sucesso, tornou possível uma suite de 49 ficheiros de teste em que a esmagadora maioria corre em memória pura, sem infraestrutura --- uma propriedade que se revelou mais valiosa em produtividade de desenvolvimento do que qualquer otimização de desempenho isolada. A reutilização semântica de ramos adaptativos por similaridade coseno em pgvector é o contributo mais original do sistema: ao decidir entre reutilização total, parcial ou geração de novo conteúdo consoante a proximidade a contextos de dificuldade já resolvidos, o custo marginal de servir cada novo Learner tende a diminuir à medida que a base de conhecimento cresce, em vez de escalar linearmente com o número de utilizadores. E a defesa em profundidade --- desde a validação de SSRF na ingestão até à sanitização de saída da IA e ao isolamento tipado de mensagens contra \textit{prompt injection} --- tratou a superfície de ataque introduzida pela IA generativa como um cidadão de primeira classe do modelo de ameaça, e não como uma reflexão tardia.

O processo de desenvolvimento também deixou lições que vale a pena registar. A auditoria conduzida para este relatório encontrou duas regras de \textit{rate limiting} inoperantes precisamente nos dois \textit{endpoints} mais utilizados da aplicação --- um lembrete de que cobertura de configuração não é o mesmo que cobertura de teste de integração, e de que a auditoria retrospetiva de um sistema em produção encontra sempre algo que a implementação inicial não previu. De forma semelhante, a avaliação de usabilidade com utilizadores reais revelou uma \textit{race condition} visual que nenhum teste automatizado tinha capturado, e expôs uma tensão de desenho --- a inscrição simultânea em múltiplos tópicos dilui o modelo de ``um desafio por dia'' que motivou a escolha do canal móvel --- que só se torna visível quando o sistema é usado, não quando é apenas especificado. Estas descobertas não diminuem o resultado; são, precisamente, o que distingue uma validação técnica de uma validação superficial.

A limitação mais honesta do projeto é também a mais simples de enunciar: o sistema foi validado com doze participantes num ambiente controlado, não com utilização real em escala. As métricas instrumentadas --- dezoito indicadores Micrometer e oito regras de alerta Prometheus --- garantem que o sistema está preparado para ser observado em produção, mas não substituem os dados que só a produção pode gerar. Latências reais sob carga concorrente, a distribuição efetiva das três estratégias de reutilização adaptativa, e a taxa de deduplicação semântica sobre um catálogo de conteúdo muito maior permanecem, nesta fase, hipóteses bem fundamentadas em vez de factos medidos.

As direções de trabalho futuro mais relevantes --- funcionalidades sociais para reforçar a retenção diária, reconhecimento formal via Open Badges e créditos ECTS na rede U!REKA, e validação automática do alinhamento de cada tarefa gerada com a competência DigComp~2.2 pretendida --- estão detalhadas no Capítulo~\ref{cap:discussao} (§\ref{sec:trabalho-futuro}) e não se repetem aqui. Fica, no entanto, a síntese da decisão estratégica que as enquadra a todas: o Play4Change pode aprofundar-se junto do público universitário tecnologicamente confortável, através do reconhecimento formal de competências, ou pode voltar-se para o perfil de menor afinidade digital que mais beneficiaria do reforço dessas competências, através de mecanismos de acompanhamento humano que reduzam a dependência exclusiva da auto-motivação. As duas direções não se excluem, mas exigem prioridades de engenharia distintas; a escolha entre elas define o que o Play4Change deve ser na sua próxima iteração.

O que este projeto demonstra, acima de tudo, é que é tecnicamente viável construir uma plataforma de aprendizagem adaptativa que gera e adapta o seu próprio conteúdo pedagógico a partir de conhecimento fornecido por um formador, sem depender de um catálogo pré-construído --- e que fazê-lo com garantias de segurança e qualidade verificáveis não é incompatível com um modelo de utilização simples, de um desafio por dia, pensado para caber na rotina de qualquer cidadão.
